\documentclass[11pt]{article}
\usepackage[margin=2.2cm]{geometry}
\usepackage{listings}
\usepackage{xcolor}
\usepackage{booktabs}
\usepackage{longtable}
\usepackage{amsmath,amsthm}
\usepackage{enumitem}
\usepackage{hyperref}
\usepackage{graphicx}
\hypersetup{colorlinks=true, linkcolor=black, urlcolor=blue}

% ----------------------------------------------------------
% SQL listing style
% ----------------------------------------------------------
\definecolor{kw}{RGB}{127,0,85}
\definecolor{str}{RGB}{26,98,26}
\definecolor{cmt}{RGB}{96,96,96}
\lstdefinelanguage{mysql}{
  keywords={SELECT,INSERT,UPDATE,DELETE,FROM,WHERE,JOIN,LEFT,RIGHT,INNER,OUTER,
            ON,GROUP,BY,ORDER,ASC,DESC,HAVING,LIMIT,CREATE,TABLE,DATABASE,DROP,
            PRIMARY,FOREIGN,KEY,REFERENCES,NOT,NULL,UNIQUE,CHECK,ENUM,DEFAULT,
            CONSTRAINT,INDEX,AFTER,BEFORE,FOR,EACH,ROW,BEGIN,END,DELIMITER,
            PROCEDURE,FUNCTION,RETURNS,DECLARE,IF,THEN,SIGNAL,SQLSTATE,SET,
            MESSAGE_TEXT,START,TRANSACTION,COMMIT,ROLLBACK,EXIT,HANDLER,FOR,
            TRIGGER,CASE,AVG,SUM,COUNT,MIN,MAX,INTERSECT,UNION,EXCEPT,ALL,
            OVER,RANK,PERCENT_RANK,VIEW,AS,IN,IS,AND,OR,CASCADE,RESTRICT,
            TIMESTAMP DIAGNOSTICS DATE TIME WEEKDAY CURDATE AUTO INCREMENT,
            BIGINT UNSIGNED INT UNSIGNED INT DECIMAL VARCHAR SMALLINT TINYINT},
  keywordstyle=\color{kw}\bfseries,
  sensitive=true,
  comment=[l]{--},
  morecomment=[s]{/*}{*/},
  commentstyle=\color{cmt}\itshape,
  string=[b]',
  stringstyle=\color{str},
  showstringspaces=false,
  basicstyle=\ttfamily\small,
  breaklines=true,
  frame=single,
  framerule=0pt,
  rulecolor=\color{gray!30},
  backgroundcolor=\color{gray!5}
}
\lstset{language=mysql}

\title{Hospital Appointment System\\
\large RDBMS Concepts Report --- RDBMS Course Project, Final Submission}
\author{Prakhar Upadhyay \and Shaurya \and Tanay}
\date{October 2026}

\begin{document}
\maketitle

\begin{abstract}
This report documents the database design of a hospital appointment
management system built on MySQL 8 with a Next.js application front end.
We present the requirement analysis, an ER model, a 1NF--3NF normalization
walkthrough, the physical schema with integrity constraints, triggers,
stored procedures and views, transaction design for concurrent booking and
billing, and the indexing strategy with EXPLAIN evidence. All SQL is shipped
in the repository as executable scripts.
\end{abstract}

\tableofcontents

% ==========================================================
\section{Introduction and Scope}
% ==========================================================
The system manages the day-to-day operations of a mid-size clinic: patient
registration, doctor schedules, appointment booking, prescriptions,
inventory-driven billing and payments, plus management reports.
Three durable artefacts document the design: \texttt{schema.sql}
(schema + sample data), \texttt{report\_queries.sql} (the analytical
queries used for screenshots), and this report.

Technology: MySQL 8, the \texttt{mysql2} driver, Next.js App Router
(server actions), Bun as runtime and package manager.

% ==========================================================
\section{Database Design}
% ==========================================================

\subsection{ER Model}
Figure~\ref{fig:er} shows the ER model. The full diagram, including both a
Mermaid rendering and a dbdiagram.io specification, lives in
\texttt{docs/er-diagram.md}.

\begin{table}[h]\centering
\begin{tabular}{@{}lll@{}}
\toprule
Relationship & Cardinality & Enforcement \\
\midrule
departments $\rightarrow$ doctors & 1:N & \texttt{doctors.dept\_id} FK \\
doctors $\rightarrow$ doctor\_schedules & 1:N & FK, \texttt{ON DELETE CASCADE} \\
patients $\rightarrow$ appointments & 1:N & FK, \texttt{ON DELETE RESTRICT} \\
doctors $\rightarrow$ appointments & 1:N & FK, \texttt{ON DELETE RESTRICT} \\
appointments $\rightarrow$ prescriptions & 1:1 & unique FK column \\
appointments $\rightarrow$ billing & 1:1 & unique FK column \\
prescriptions $\leftrightarrow$ medicines & M:N & bridge \texttt{prescription\_items} \\
billing $\rightarrow$ payments & 1:N & FK, \texttt{ON DELETE RESTRICT} \\
users $\rightarrow$ doctors & 0..1 & nullable FK, \texttt{SET NULL} \\
\bottomrule
\end{tabular}
\caption{Cardinality matrix for the schema.}\label{fig:er}
\end{table}

\subsection{Normalization Walkthrough}

\paragraph{Unnormalized form.} A visit sheet records
(patient, doctor, date, time, fee, drug 1..3, doses, days, bill status,
payment method, amount). Repeating-degree groups such as
\texttt{medicine1..medicine3} violate 1NF.

\paragraph{1NF.} Atomic values only; the repeating group becomes one row per
prescribed medicine, which eventually forms \texttt{prescription\_items}.

\paragraph{2NF.} Assuming the composite line key
\texttt{(patient, doctor, date, time, medicine)}, several attributes depend on
only part of the key: \texttt{consultation\_fee} depends on the doctor part;
\texttt{dob}, \texttt{phone} depend on the patient part; \texttt{dosage},
\texttt{days} depend on the medicine line. These partial dependencies are
removed by splitting out \texttt{patients}, \texttt{doctors},
\texttt{appointments} and the prescription line table.

\paragraph{3NF.} Transitive chains inside appointment rows
($\text{appointment} \rightarrow \text{doctor} \rightarrow \text{department}$
and $\text{appointment} \rightarrow \text{prescription} \rightarrow \text{medicines}$)
are removed by the final 3NF schema shown in Section~\ref{sec:physical}.
The only retained derived attribute is \texttt{billing.total}, an intentional
financial snapshot: a bill must not silently reprice when a medicine price
changes later, so \texttt{total} is stored with correct values maintained by
the \texttt{generate\_bill} procedure.

% ==========================================================
\section{Physical Schema}\label{sec:physical}
% ==========================================================

The design enforces integrity \emph{physically}, not by convention:

\begin{itemize}
  \item Primary keys on all 12 tables (surrogate \texttt{AUTO\_INCREMENT}).
  \item Foreign keys with explicit delete behaviour:
        \texttt{RESTRICT} prevents deleting a doctor or patient with
        appointments; \texttt{CASCADE} keeps schedules, prescriptions and
        bills attached to their owner; \texttt{SET NULL} detaches a doctor
        login from a deleted doctor.
  \item \texttt{UNIQUE (doctor\_id, appt\_date, appt\_time)} --- a single
        appointment per doctor per hour slot, the physical double-booking
        guard.
  \item \texttt{CHECK} constraints:
        \texttt{consultation\_fee > 0}, \texttt{stock >= 0},
        \texttt{total >= 0}, \texttt{amount > 0},
        \texttt{start\_time < end\_time}, \texttt{dob <= CURDATE()}.
  \item \texttt{ENUM} types constrain \texttt{status},
        \texttt{gender}, \texttt{method}, \texttt{role} to their
        nominal domains, with \texttt{DEFAULT 'booked'} as the
        appointment status.
\end{itemize}

\begin{lstlisting}
CREATE TABLE appointments (
  appointment_id INT UNSIGNED AUTO_INCREMENT PRIMARY KEY,
  patient_id INT UNSIGNED NOT NULL,
  doctor_id  INT UNSIGNED NOT NULL,
  appt_date  DATE NOT NULL,
  appt_time  TIME NOT NULL,
  status ENUM('booked','completed','cancelled') NOT NULL DEFAULT 'booked',
  reason VARCHAR(255) NULL,
  created_at TIMESTAMP DEFAULT CURRENT_TIMESTAMP,
  UNIQUE KEY uq_doctor_slot (doctor_id, appt_date, appt_time),
  CONSTRAINT fk_appt_patient FOREIGN KEY (patient_id)
    REFERENCES patients(patient_id) ON DELETE RESTRICT ON UPDATE CASCADE,
  CONSTRAINT fk_appt_doctor FOREIGN KEY (doctor_id)
    REFERENCES doctors(doctor_id) ON DELETE RESTRICT ON UPDATE CASCADE
);
\end{lstlisting}

% ==========================================================
\section{Triggers}
% ==========================================================
Triggers encode business rules atomically at the row level.
The sealed fiscal rules keep stock, billing and audit consistent without
relying on every caller remembering three update statements.

\begin{enumerate}
  \item \textbf{\texttt{trg\_item\_insert}} (AFTER INSERT on
        \texttt{prescription\_items}) --- decrement the medicine stock by one
        for every written prescription line; guarded by \texttt{stock > 0}
        so stock never dips below the \texttt{CHECK} limit.
  \item \textbf{\texttt{trg\_appt\_completed}} (AFTER UPDATE on
        \texttt{appointments}) --- when an appointment transitions to
        \texttt{'completed'}, insert the initial billing row using the
        doctor's consultation fee.
  \item \textbf{\texttt{trg\_payment\_insert}} (AFTER INSERT on
        \texttt{payments}) --- recompute the cumulative paid amount and
        flip \texttt{billing.paid\_status} to \texttt{'paid'} once the
        total is covered.
  \item \textbf{\texttt{trg\_appt\_audit\_update}} and
        \textbf{\texttt{trg\_appt\_audit\_delete}} --- insert rows in the
        \texttt{audit\_log} for every appointment modification.
\end{enumerate}

\begin{lstlisting}
CREATE TRIGGER trg_payment_insert
AFTER INSERT ON payments
FOR EACH ROW
BEGIN
  DECLARE paid_sum DECIMAL(10,2);
  SELECT COALESCE(SUM(amount), 0) INTO paid_sum
    FROM payments WHERE bill_id = NEW.bill_id;
  UPDATE billing
     SET paid_status = IF(paid_sum >= total, 'paid', 'unpaid')
   WHERE bill_id = NEW.bill_id;
END
\end{lstlisting}

% ==========================================================
\section{Stored Procedures and Functions}
% ==========================================================
\begin{itemize}
  \item \texttt{book\_appointment(patient, doctor, date, time, reason)} ---
        validates that the slot lies within the doctor's schedule
        (WEEKDAY + 1 maps a date to the day-of-week encoding
        $1{=}\text{Mon} \dots 7{=}\text{Sun}$), opens a transaction, performs
        the insert and commits; a SQL exception rolls back and
        \texttt{RESIGNAL}s the original error to the caller.
  \item \texttt{cancel\_appointment(id)} --- sets the status to
        \texttt{'cancelled'}, deletes any recorded payments (refund) and the
        generated bill row; refuses to cancel twice.
  \item \texttt{generate\_bill(appointment\_id)} --- accumulates the
        consultation fee plus the price sum of prescribed medicines and
        refreshes the bill.
  \item \texttt{get\_available\_slots(doctor, date)} --- reports schedule
        ranges and the count of occupied 30-minute cells.
  \item \texttt{calculate\_age(dob)} \emph{(function)} ---
        \texttt{TIMESTAMPDIFF(YEAR, dob, CURDATE())}; used by views and
        reports for age grouping.
\end{itemize}

% ==========================================================
\section{Transactions (ACID)}
% ==========================================================
Booking is the race-sensitive write: two receptionists can hit the same
doctor/date at once. The insert path (server action
\texttt{bookAppointmentAction}) wraps three statements in one transaction:

\begin{lstlisting}
START TRANSACTION;
  -- 1. validate the doctor's schedule for that day
  SELECT 1 FROM doctor_schedules
   WHERE doctor_id = ? AND day_of_week = WEEKDAY(?) + 1
     AND start_time <= ? AND end_time > ? LIMIT 1;
  -- 2. intention-lock every row of this doctor/date
  SELECT appointment_id FROM appointments
   WHERE doctor_id = ? AND appt_date = ? FOR UPDATE;
  -- 3. insert; a concurrent insert blocks here and,
  --    when it resumes, violates uq_doctor_slot and
  --    ER_DUP_ENTRY surfaces as a friendly error
  INSERT INTO appointments
    (patient_id, doctor_id, appt_date, appt_time, status, reason)
  VALUES (?, ?, ?, ?, 'booked', NULLIF(?, ''));
COMMIT;
\end{lstlisting}

\textbf{Atomicity:} all-or-nothing insert. \textbf{Consistency:} schedule and
uniqueness invariants survive every concurrent schedule. \textbf{Isolation:}
\texttt{FOR UPDATE} at \texttt{REPEATABLE READ} serialises competing booking
transactions on the shared predicate. \textbf{Durability:} the \texttt{COMMIT}
plus InnoDB redo log.

The billing path is equally transactional: insert \texttt{payments}, verify
the cumulative remaining balance afterwards and roll back on overpay errors;
the trigger keeps \texttt{paid\_status} consistent within the same commit.

% ==========================================================
\section{Views}
% ==========================================================
\begin{itemize}
  \item \texttt{todays\_appointments} --- patient, doctor, department, time
        and status for \texttt{CURDATE()}.
  \item \texttt{doctor\_workload} --- total appointments per doctor including
        zero-load doctors (\texttt{LEFT JOIN}).
  \item \texttt{pending\_bills} --- unpaid bills with patient and doctor names.
  \item \texttt{low\_stock\_medicines} --- stock below the restock threshold
        of 10, driving the dashboard widget and warning badges.
\end{itemize}

% ==========================================================
\section{Analytical Requirements}
% ==========================================================
The report screenshots map one-to-one to \texttt{report\_queries.sql}:

\begin{itemize}
  \item \textbf{JOIN} of appointments with patient, doctor and department.
  \item \textbf{GROUP BY/HAVING}: departments with more than 10 appointments
        this month.
  \item \textbf{Subquery}: doctors with above-average workload via
        $\texttt{total} > (\text{AVG of totals})$.
  \item \textbf{Aggregate}: revenue per department per month.
  \item \textbf{Window function}: \texttt{RANK() OVER (ORDER BY
        total\_appointments DESC)} plus \texttt{PERCENT\_RANK}.
  \item \textbf{Set operations}: patients treated in both General Medicine
        and Cardiology \texttt{INTERSECT}ed and \texttt{UNION}ed by month.
\end{itemize}

% ==========================================================
\section{Indexing Strategy}
% ==========================================================
The booking path filters on \texttt{(doctor\_id, appt\_date)}: the
composite index serves both the unique constraint and the read, so lookups
for \texttt{EXPLAIN} on that predicate change from scanning all rows to a
\texttt{ref}/{range} plan. Phone recovery in the front desk search is the
first read path; it relies on a seeded phone search demo (both EXPLAIN
before/after runs proved non-unique key scan/range choice in Workbench,
both recorded for the report).

A second relevant detail: the \texttt{uq\_doctor\_slot} constraint itself is
an index, so the FOR UPDATE lock predicate and the uniqueness check share
one lookup path per booking.

% ==========================================================
\section{Application Architecture}
% ==========================================================
Next.js App Router splits the system as follows:

\begin{itemize}
  \item \textbf{Server actions} (\texttt{src/server/actions.ts}) validate
        inputs with Zod and issue parameterised SQL. No user input is
        concatenated, blocking SQL injection at the source.
  \item \textbf{Authentication} uses JWTs via \texttt{jose} with passwords
        \texttt{bcrypt}-verified from \texttt{users}; the role decorates the
        session and the doctor login optionally binds to \texttt{doctor\_id}.
  \item \textbf{Pages} are server components fetching with the
        \texttt{mysql2/promise} pool; the interactive book form is a client
        component driven by URL search params, so every step is linkable.
  \item \textbf{Dashboard} is a draggable widget grid where each tile is a
        live query result (appointments, revenue, bills, stock, workload).
\end{itemize}

% ==========================================================
\section{Conclusion}
% ==========================================================
The design pushes invariants down into the database layer (unique slots,
checks, triggers for stock and billing, procedures for workflows) while the
application concentrates on presentation and transaction orchestration.
Deliverables --- the ER model, the normalization walkthrough, the
\dots{}complete \texttt{schema.sql}, the analytical query set and this
report --- reproduce every artefact used for the final demonstration.

\vspace{1em}
\begin{center}\rule{8cm}{0.4pt}\end{center}

\begin{center}
\emph{Built by Prakhar, Shaurya and Tanay --- RDBMS Project, Final.}
\end{center}

\end{document}
